\documentclass[11pt,a4paper]{article}
\newcommand{\Rodzaj}{Szablon sprawozdania}
\usepackage{fontspec}
\setmainfont{texgyretermes-regular.otf}[BoldFont=texgyretermes-bold.otf,ItalicFont=texgyretermes-italic.otf]
\usepackage[polish]{babel}
\usepackage{amsmath,amssymb,geometry,array,tabularx,graphicx,caption,fancyhdr}
\usepackage[table]{xcolor}
\geometry{left=18mm,right=18mm,top=19mm,bottom=19mm,headheight=14pt,headsep=6mm}
\setlength{\parindent}{0pt}\setlength{\parskip}{5pt}
\setlength{\emergencystretch}{2em}
\renewcommand{\arraystretch}{1.3}\setlength{\tabcolsep}{5pt}
\setlength{\fboxsep}{6pt}
\setlength{\abovedisplayskip}{8pt}\setlength{\belowdisplayskip}{8pt}
\setlength{\abovedisplayshortskip}{5pt}\setlength{\belowdisplayshortskip}{5pt}
\widowpenalty=10000\clubpenalty=10000
\raggedbottom
\newcolumntype{Y}{>{\raggedright\arraybackslash}X}
\newcolumntype{C}{>{\centering\arraybackslash}X}
\newcommand{\naglowek}{\rowcolor{gray!12}}
\newcommand{\odstep}{\par\smallskip}
\captionsetup{font=small,labelfont=bf,justification=justified,singlelinecheck=false}
\pagestyle{fancy}\fancyhf{}\fancyhead[L]{\small Ćwiczenie 83 — spektrometr optyczny}
\fancyhead[R]{\small\Rodzaj}\fancyfoot[C]{\small\thepage}
\renewcommand{\headrulewidth}{0.3pt}
\makeatletter
\renewcommand\section{\@startsection{section}{1}{0pt}{10pt}{4pt}{\large\bfseries}}
\renewcommand\subsection{\@startsection{subsection}{2}{0pt}{8pt}{4pt}{\normalsize\bfseries}}
\makeatother
\newcommand{\pole}{\rule{23mm}{0.3pt}}
\newcommand{\miejsce}[2]{\par\medskip\noindent\fbox{\begin{minipage}[t][#1][t]{\dimexpr\linewidth-2\fboxsep-2\fboxrule\relax}\small #2\end{minipage}}\par\medskip}
\newcommand{\metryka}{%
\begingroup\small\setlength{\parskip}{0pt}\renewcommand{\arraystretch}{1.1}
\begin{tabularx}{\linewidth}{|p{29mm}|Y|p{18mm}|p{18mm}|p{19mm}|}\hline
\textbf{Wydział}\newline Informatyki & \textbf{Imię i nazwisko}\newline 1. Marcin Podyma\newline 2. Paweł Mierzwa & \textbf{Rok}\newline 2026 & \textbf{Grupa}\newline 3 & \textbf{Zespół}\newline 1\\\hline
\end{tabularx}\par\nointerlineskip
\begin{tabularx}{\linewidth}{|p{29mm}|Y|p{19mm}|}
\parbox[t]{29mm}{\centering\textbf{PRACOWNIA\\FIZYCZNA\\WFiIS AGH}} & \textbf{Temat:}\newline Spektrometr optyczny — widmo liniowe par rtęci & \textbf{Nr}\newline\textbf{ćwiczenia}\newline 83\\\hline
\end{tabularx}\par\nointerlineskip
\begin{tabularx}{\linewidth}{|p{29mm}|Y|Y|Y|Y|p{19mm}|}
\textbf{Data wykonania}\newline\strut & \textbf{Data oddania}\newline\strut & \textbf{Zwrot do popr.} & \textbf{Data oddania} & \textbf{Data zaliczenia} & \textbf{OCENA}\\\hline
\end{tabularx}\endgroup\medskip}

\begin{document}
\metryka
\section{Cel ćwiczenia}
Wyznaczenie długości fal widma liniowego par rtęci za pomocą spektrometru
z siatką dyfrakcyjną oraz identyfikacja odpowiadających im przejść energetycznych.
\section{Podstawy teoretyczne}
Widmo liniowe wynika z przejść między dyskretnymi poziomami energetycznymi.
Energia emitowanego fotonu spełnia zależność
\[
\boxed{E_g-E_d=h\nu=\frac{hc}{\lambda}.}
\]
Dla światła padającego prostopadle na siatkę warunek maksimum ma postać
\[
\boxed{d\sin\alpha_n=n\lambda,\qquad
\lambda_n=\frac{d\sin\alpha_n}{n},\qquad d=\frac1N.}
\]
Tutaj $d$ jest stałą siatki, $N$ liczbą rys na jednostkę długości,
$n=1,2$ dodatnim numerem rzędu, a $\alpha_n$ kątem odchylenia od
wiązki nieodchylonej. Rząd $n=0$ służy wyznaczeniu kierunku zerowego,
nie oblicza się z niego długości fali.
\section{Aparatura i przebieg pomiarów}
Stanowisko obejmuje lampę rtęciową, kolimator z regulowaną szczeliną,
siatkę na stoliku oraz obracaną lunetę z podziałką kątową i noniuszem.
\begin{center}\small\fbox{Lampa Hg $\longrightarrow$ kolimator
$\longrightarrow$ siatka $\longrightarrow$ luneta i odczyt kąta}\end{center}
\begin{tabularx}{\linewidth}{|Y|Y|}\hline
\naglowek\textbf{Parametr} & \textbf{Wartość rzeczywista}\\\hline
Stała siatki $d$ / liczba rys $N$ & \rule{0pt}{7mm}\\\hline
Szerokość szczeliny & \rule{0pt}{7mm}\\\hline
Przyjęta dokładność odczytu kąta & \rule{0pt}{7mm}\\\hline
Kierunek zerowy $\theta_0$ & \rule{0pt}{7mm}\\\hline
\end{tabularx}
\miejsce{19mm}{Opis rzeczywistego przebiegu: ustawienie i kontrola siatki,
odczyt zera, pomiar obu stron w I i II rzędzie, liczba powtórzeń,
osobne odczyty dwóch linii żółtych.}

\clearpage
\section{Wyniki pomiarów}
Zestaw odchylenia lewej i prawej strony dla każdej linii i obu rzędów.
Surowe wskazania skali $\theta$ oraz odczyt zera $\theta_0$ zachowaj
na karcie pomiarowej. Przy powtórzeniach podaj ich liczbę
i sposób uśredniania. Identyfikator linii musi być taki sam w obu rzędach.
\begin{center}\small
\captionof{table}{Odchylenia obu stron wyznaczone z własnych odczytów.}
\begin{tabularx}{\linewidth}{|c|c|Y|C|C|}\hline
\naglowek\textbf{ID} & \textbf{Rząd} & \textbf{Barwa} &
$\alpha_L$ [$^\circ$] & $\alpha_P$ [$^\circ$]\\\hline
\csname @@input\endcsname{odchylenia_szablon.tex}
\end{tabularx}\end{center}
\miejsce{35mm}{Kontrola ustawienia: różnica odchyleń stron w I rzędzie,
porównanie z progiem około $10'$ oraz rzeczywiste poprawki ustawienia.
Nie stwierdzaj eliminacji wszystkich błędów systematycznych tylko na podstawie uśrednienia stron.}

\clearpage
\section{Opracowanie kątów i długości fal}
\subsection{Kąt odchylenia}
Odczyt $\theta$ jest wskazaniem skali, której zero nie musi odpowiadać
wiązce nieodchylonej. Stopnie i minuty przeliczamy jako
\[
\boxed{\theta= D+\frac{M}{60}\quad[{}^\circ],\qquad
\theta_{\mathrm{rad}}=\theta_{\mathrm{deg}}\frac{\pi}{180}.}
\]
Różnicę względem $\theta_0$ należy sprowadzić do przedziału
$(-180^\circ,180^\circ]$, aby poprawnie przejść przez zero skali.
Dla pary lewego i prawego prążka obliczamy
\[
\boxed{\begin{gathered}
\alpha_L=|\operatorname{wrap}(\theta_L-\theta_0)|,\qquad
\alpha_P=|\operatorname{wrap}(\theta_P-\theta_0)|,\\[3pt]
\overline\alpha_n=\frac{\alpha_L+\alpha_P}{2},\qquad
\lambda_n=\frac{d\sin\overline\alpha_n}{n}.
\end{gathered}}
\]
Najpierw wyznaczamy średni kąt, następnie długość fali.
W obliczeniach funkcji sinus należy użyć właściwej jednostki kąta.
Dla powtórzeń najpierw uśrednia się odchylenia w obrębie tej samej linii i rzędu.
\subsection{Zestawienie wyników obu rzędów}
\begin{center}\small
\captionof{table}{Kąty i długości fal tych samych linii w I i II rzędzie.}
\begin{tabularx}{\linewidth}{|c|Y|C|C|C|C|C|}\hline
\naglowek\textbf{ID} & \textbf{Barwa} & $\overline\alpha_1$ [$^\circ$] &
$\lambda_1$ [nm] & $\overline\alpha_2$ [$^\circ$] &
$\lambda_2$ [nm] & $|\lambda_1-\lambda_2|$ [nm]\\\hline
\csname @@input\endcsname{wyniki_rzedy.tex}
\end{tabularx}\end{center}
\miejsce{37mm}{Przykład obliczenia dla jednej linii: surowe odczyty,
przejście przez zero skali, średni kąt, stała siatki i długość fali z jednostkami.}

\clearpage
\section{Identyfikacja linii widmowych}
Dla każdej linii zestaw wynik z długością fali odpowiadającego jej
przejścia na diagramie poziomów Hg w instrukcji. Obie linie dubletu
żółtego wpisz w osobnych wierszach. Symbole poziomów zapisuj
w postaci $N^{2S+1}X_J$, z poziomem początkowym o wyższej energii.
\[
\boxed{\lambda_{\mathrm{konc}}=\frac{\lambda_1+\lambda_2}{2}
\quad\text{(dla linii zmierzonej w obu rzędach)},\qquad
\Delta E=\frac{hc}{\lambda_{\mathrm{konc}}}.}
\]
Średnia arytmetyczna jest przyjętym sposobem zestawienia dwóch rzędów.
Gdy pomiar istnieje tylko w jednym rzędzie, wpisz ten wynik i zaznacz brak
drugiego. Nie obliczaj wtedy różnicy między rzędami dla tej linii.
\begin{center}\small
\captionof{table}{Wyniki końcowe i identyfikacja przejść energetycznych.}
\begin{tabularx}{\linewidth}{|c|Y|C|C|Y|Y|}\hline
\naglowek\textbf{ID} & \textbf{Barwa} & $\lambda_{\mathrm{konc}}$ [nm] &
$\lambda_{\mathrm{ref}}$ [nm] & \textbf{Poziom początkowy} & \textbf{Poziom końcowy}\\\hline
\csname @@input\endcsname{identyfikacja.tex}
\end{tabularx}\end{center}
\miejsce{30mm}{Uzasadnienie przypisania linii: zgodność długości fal,
rozdzielenie dubletu i ewentualne niejednoznaczności.}
\miejsce{48mm}{Opcjonalny schemat zidentyfikowanych przejść lub diagram linii widmowych.
Instrukcja wymaga tabeli identyfikacji; wykres nie jest obowiązkowy.}

\clearpage
\section{Ocena błędu pomiaru}
Instrukcja wskazuje, że istotne mogą być błędy systematyczne związane
z niedoskonałym ustawieniem spektrometru. Jako oszacowanie błędu
zaleca średnią różnicę długości fal uzyskanych z I i II rzędu.
Dla $K$ linii zmierzonych w obu rzędach przyjmujemy dodatnią miarę rozbieżności:
\[
\boxed{\Delta\lambda_{\mathrm{szac}}=
\frac1K\sum_{i=1}^{K}|\lambda_{i,1}-\lambda_{i,2}|.}
\]
Jest to szacunkowa miara błędu według instrukcji, a nie automatycznie
niepewność standardowa ani rozszerzona. Podaj, które linie uwzględniono.
\miejsce{22mm}{Oszacowanie błędu: liczba wspólnych linii $K$, różnice między
rzędami i obliczona wartość $\Delta\lambda_{\mathrm{szac}}$.}
Jeżeli dostępne są niepewności stałej siatki i kąta, dodatkowo można
zastosować propagację dla niezależnych wielkości:
\[
\boxed{u(\lambda_n)=\frac1n\sqrt{
\sin^2\alpha_n\,u^2(d)+d^2\cos^2\alpha_n\,u^2(\alpha_n)}.}
\]
W tym wzorze $u(\alpha_n)$ musi być wyrażona w radianach.
Sam zalecany odczyt do $1'$ nie jest gotowym bilansem niepewności.
\miejsce{17mm}{Dyskusja wyników: różnice obu stron, ewentualny wspólny kierunek
różnic I--II rzędu, wpływ ustawienia siatki i identyfikacji prążków.
Odróżnij zaobserwowane rozbieżności od przypuszczalnych przyczyn.}
\section{Wnioski}
\begin{tabularx}{\linewidth}{|>{\raggedright\arraybackslash}p{43mm}|Y|}\hline
\naglowek\textbf{Zagadnienie} & \textbf{Wynik i interpretacja}\\\hline
Wyznaczone długości fal & \rule{0pt}{13mm}\\\hline
Zgodność I i II rzędu oraz ocena błędu & \rule{0pt}{13mm}\\\hline
Identyfikacja przejść i dublet żółty & \rule{0pt}{13mm}\\\hline
Ograniczenia pomiaru & \rule{0pt}{13mm}\\\hline
\end{tabularx}
\end{document}
